\documentclass[11pt]{article}
\usepackage{amsmath,amssymb,geometry,hyperref}
\geometry{margin=1in}
\title{Replication of the Theory Core of\\ ``Observation of Moir\'e Plasmonic Skyrmion Clusters''\\ (Zhang et al., arXiv:2411.05576)}
\author{Automated replication (textures-100 pipeline)}
\date{}
\begin{document}
\maketitle

\section*{Recommendation: REPLICATE (partial). Verdict: PASS-PARTIAL.}
\textbf{C/A scores:} Correctness (C) $=0.6$, Alignment (A) $=0.7$.
The paper is \emph{not} an experiment-only drop: it carries a self-contained,
reproducible analytic theory core (an SPP interference model plus a
skyrmion-number-density integral) alongside its FDTD simulations and near-field
experiment. We reproduced the analytic core from scratch.

\section{What the paper is}
Zhang et al.\ report ``moir\'e plasmonic skyrmion clusters'': two twisted
hexagonal groups of nanoslits etched in a gold film launch surface plasmon
polaritons (SPPs) whose interference sculpts the axial ($z$) electric field into
optical-skyrmion superlattices. Twisting one group by $\theta$ produces
crystallized (commensurate) and quasi-crystallized (incommensurate) skyrmion
lattices, and rapidly inverts the per-cell skyrmion number. Evidence: 3D FDTD
(Maxwell) simulations $+$ near-field microscopy of fabricated samples.

\section{Reproducible theory core (extracted, coded from scratch)}
\begin{align}
E_z(\mathbf r) &= e^{-|k_z|z}\Big\{\textstyle\sum_{j=1}^{6}\tilde E_j
   e^{-i(\mathbf k_j(0)\cdot\mathbf r-\phi_j)}
 + \sum_{j=1}^{6}\tilde E_j e^{-i(\mathbf k_j(\theta)\cdot\mathbf r-\phi_j-\Delta\phi)}\Big\}, \tag{1}\\
\mathbf k_j(\theta) &= \mathbb{R}(\theta)\,k_t[\cos\varphi_j,\sin\varphi_j],\quad
   k_t^2+k_z^2=k_0^2,\quad \phi_j=2\pi\sigma_j/N,\\
\bar{\mathbf E}(\mathbf r,\theta) &=
   \frac{\mathrm{Re}[\mathbf E_a(\mathbf r,0)+\mathbf E_b(\mathbf r,\theta)]}
        {|\,\mathrm{Re}[\dots]\,|}, \tag{2}\\
Q(\theta) &= \frac{1}{4\pi}\iint_{\mathrm{cell}} s(\mathbf r,\theta)\,d^2r, \qquad
s = \bar{\mathbf E}\cdot(\partial_x\bar{\mathbf E}\times\partial_y\bar{\mathbf E}). \tag{3,4}
\end{align}
The transverse field is fixed per partial wave from $\nabla\cdot\mathbf E=0$:
$\mathbf E_t=-i(|k_z|/k_t)E_z\,\hat{\mathbf k}_j$. We evaluate $Q$ with the
Berg--L\"uscher lattice solid-angle method (integer-robust) and a
finite-difference cross-check.

\section{Results}
\begin{center}
\begin{tabular}{lccc}
\hline
Case & Paper claim & This work (Berg) & Match \\ \hline
Elementary single hex SPP lattice & $Q=\pm1$ per cell & $Q=+1.00$ per cell & \textbf{yes} \\
Composite twist $\theta=38.21^\circ$ & $Q=-3$ per moir\'e unit & $Q\approx0$ total; $\pm$ nucleation & partial \\
Twist sweep & rapid $Q$ inversion & balanced $\pm$ skyrmions & qualitative \\ \hline
\end{tabular}
\end{center}

The \textbf{elementary lattice is reproduced exactly} ($Q=+1$ per primitive
cell), which validates the full Eq.(1)$\to$(4) $+$ Berg--L\"uscher pipeline. The
specific cluster value $Q=-3$ per moir\'e unit is \emph{not} isolated using
main-text parameters alone: over the symmetric composite grid the total charge
is $\approx0$, i.e.\ equal numbers of $+$ and $-$ skyrmions nucleate. This is
consistent with the paper's red/blue (positive/negative) topological-point
picture, but the exact $-3$ requires the Supplemental Material definitions
(nanoslit index map $\sigma_j$, the precise moir\'e primitive cell over which the
integral is restricted, and the S5 critical-inversion condition) that are not in
the extracted main text. Hence \emph{partial} correctness.

\section*{Provenance}
Berg--L\"uscher solid-angle kernel adapted from
\texttt{\char`~/shared-kernels-cache/ollie\_berg\_luscher\_topological\_charge\_kernel.py}.
CPU-only, numpy only.
\end{document}
